\documentclass[aspectratio=169, xcolor=dvipsnames]{beamer}

% --- Theme Setup ---
\usetheme{Madrid}
\usecolortheme{default}

% --- Packages ---
\usepackage{graphicx}
\usepackage{xcolor}
\usepackage{booktabs}
\usepackage{hyperref}
\usepackage{adjustbox}
\usepackage{amssymb}
\usepackage{amsmath}
\usepackage{listings}
\usepackage{caption}

\lstset{
  basicstyle=\ttfamily\small,
  keywordstyle=\color{blue}\bfseries,
  stringstyle=\color{red},
  showstringspaces=false,
  breaklines=true
}

% --- Title Information ---
\title{Module 13}
\subtitle{Intermediate SQL/2}
\author{Milav Dabgar}
\institute{www.milav.in}
\date{\today}

\begin{document}

% Title Slide
\begin{frame}
    \titlepage
\end{frame}

\begin{frame}{Module Recap}
    \begin{itemize}
        \item Nested subquery in SQL
        \item Processes for data modification
    \end{itemize}
\end{frame}

\begin{frame}{Module Objectives}
    \begin{itemize}
        \item To learn SQL expressions for Join
        \item To learn SQL expressions for Views
    \end{itemize}
\end{frame}

\begin{frame}{Module Outline}
    \begin{itemize}
        \item Join Expressions
        \begin{itemize}
            \item Cross Join
            \item Inner Join
            \item Outer Join
            \item Left Outer Join
            \item Right Outer Join
            \item Full Outer Join
        \end{itemize}
        \item Views
        \begin{itemize}
            \item View Expansion
            \item View Update
            \item Materialized Views
        \end{itemize}
    \end{itemize}
\end{frame}

\section{Join Expressions}
\begin{frame}
  \vfill
  \centering
  \begin{beamercolorbox}[sep=8pt,center,shadow=true,rounded=true]{title}
    \usebeamerfont{title}Join Expressions\par%
  \end{beamercolorbox}
  \vfill
\end{frame}

\begin{frame}{Joined Relations}
    \begin{itemize}
        \item Join operations take two relations and return as a result another relation
        \item A join operation is a Cartesian product which requires that tuples in the two relations match (under some condition).
        \item It also specifies the attributes that are present in the result of the join
        \item The join operations are typically used as subquery expressions in the \texttt{from} clause
    \end{itemize}
\end{frame}

\begin{frame}{Types of Join between Relations}
    \begin{itemize}
        \item \textbf{Cross join}
        \item \textbf{Inner join}
        \begin{itemize}
            \item Equi-join
            \item Natural join
        \end{itemize}
        \item \textbf{Outer join}
        \begin{itemize}
            \item Left outer join
            \item Right outer join
            \item Full outer join
        \end{itemize}
        \item \textbf{Self-join}
    \end{itemize}
\end{frame}

\begin{frame}[fragile]{Cross Join}
    \begin{itemize}
        \item \texttt{CROSS JOIN} returns the Cartesian product of rows from tables in the join
        \item \textbf{Explicit}
        \begin{lstlisting}[language=SQL]
select * 
from employee cross join department;
        \end{lstlisting}
        \item \textbf{Implicit}
        \begin{lstlisting}[language=SQL]
select * 
from employee, department;
        \end{lstlisting}
    \end{itemize}
\end{frame}

\begin{frame}{Join operations - Example}
    \begin{itemize}
        \item Relation \texttt{course}
        \item Relation \texttt{prereq}
        \item Observe that prereq information is missing for CS-315 and course information is missing for CS-347
    \end{itemize}
    \begin{columns}
        \begin{column}{0.65\textwidth}
            \textbf{course}
            \begin{center}
                \begin{adjustbox}{max width=\textwidth}
                \begin{tabular}{|l|l|l|l|}
                \hline
                \textbf{course\_id} & \textbf{title} & \textbf{dept\_name} & \textbf{credits} \\ \hline
                BIO-301 & Genetics & Biology & 4 \\ \hline
                CS-190 & Game Design & Comp. Sci. & 4 \\ \hline
                CS-315 & Robotics & Comp. Sci. & 3 \\ \hline
                \end{tabular}
                \end{adjustbox}
            \end{center}
        \end{column}
        \begin{column}{0.35\textwidth}
            \textbf{prereq}
            \begin{center}
                \begin{adjustbox}{max width=\textwidth}
                \begin{tabular}{|l|l|}
                \hline
                \textbf{course\_id} & \textbf{prereq\_id} \\ \hline
                BIO-301 & BIO-101 \\ \hline
                CS-190 & CS-101 \\ \hline
                CS-347 & CS-101 \\ \hline
                \end{tabular}
                \end{adjustbox}
            \end{center}
        \end{column}
    \end{columns}
\end{frame}

\begin{frame}[fragile]{Inner Join}
    \begin{itemize}
        \item \texttt{course inner join prereq}
    \end{itemize}
    \begin{center}
        \begin{adjustbox}{max width=\textwidth}
        \begin{tabular}{|l|l|l|l|l|l|}
        \hline
        \textbf{course\_id} & \textbf{title} & \textbf{dept\_name} & \textbf{credits} & \textbf{prereq\_id} & \textbf{course\_id} \\ \hline
        BIO-301 & Genetics & Biology & 4 & BIO-101 & BIO-301 \\ \hline
        CS-190 & Game Design & Comp. Sci. & 4 & CS-101 & CS-190 \\ \hline
        \end{tabular}
        \end{adjustbox}
    \end{center}
    
    \vspace{1em}
    \begin{itemize}
        \item If specified as \texttt{natural}, the 2nd \texttt{course\_id} field is skipped
    \end{itemize}
    \begin{center}
        \begin{adjustbox}{max width=0.8\textwidth}
        \begin{tabular}{|l|l|l|l|}
        \hline
        \textbf{course\_id} & \textbf{title} & \textbf{dept\_name} & \textbf{credits} \\ \hline
        BIO-301 & Genetics & Biology & 4 \\ \hline
        CS-190 & Game Design & Comp. Sci. & 4 \\ \hline
        \end{tabular}
        \end{adjustbox}
    \end{center}
\end{frame}

\begin{frame}{Outer Join}
    \begin{itemize}
        \item An extension of the join operation that avoids loss of information
        \item Computes the join and then adds tuples from one relation that does not match tuples in the other relation to the result of the join
        \item Uses null values
    \end{itemize}
\end{frame}

\begin{frame}[fragile]{Left Outer Join}
    \begin{itemize}
        \item \texttt{course natural left outer join prereq}
    \end{itemize}
    \begin{center}
        \begin{adjustbox}{max width=0.9\textwidth}
        \begin{tabular}{|l|l|l|l|l|}
        \hline
        \textbf{course\_id} & \textbf{title} & \textbf{dept\_name} & \textbf{credits} & \textbf{prereq\_id} \\ \hline
        BIO-301 & Genetics & Biology & 4 & BIO-101 \\ \hline
        CS-190 & Game Design & Comp. Sci. & 4 & CS-101 \\ \hline
        CS-315 & Robotics & Comp. Sci. & 3 & null \\ \hline
        \end{tabular}
        \end{adjustbox}
    \end{center}
\end{frame}

\begin{frame}[fragile]{Right Outer Join}
    \begin{itemize}
        \item \texttt{course natural right outer join prereq}
    \end{itemize}
    \begin{center}
        \begin{adjustbox}{max width=0.9\textwidth}
        \begin{tabular}{|l|l|l|l|l|}
        \hline
        \textbf{course\_id} & \textbf{title} & \textbf{dept\_name} & \textbf{credits} & \textbf{prereq\_id} \\ \hline
        BIO-301 & Genetics & Biology & 4 & BIO-101 \\ \hline
        CS-190 & Game Design & Comp. Sci. & 4 & CS-101 \\ \hline
        CS-347 & null & null & null & CS-101 \\ \hline
        \end{tabular}
        \end{adjustbox}
    \end{center}
\end{frame}

\begin{frame}[fragile]{Full Outer Join}
    \begin{itemize}
        \item \texttt{course natural full outer join prereq}
    \end{itemize}
    \begin{center}
        \begin{adjustbox}{max width=0.9\textwidth}
        \begin{tabular}{|l|l|l|l|l|}
        \hline
        \textbf{course\_id} & \textbf{title} & \textbf{dept\_name} & \textbf{credits} & \textbf{prereq\_id} \\ \hline
        BIO-301 & Genetics & Biology & 4 & BIO-101 \\ \hline
        CS-190 & Game Design & Comp. Sci. & 4 & CS-101 \\ \hline
        CS-315 & Robotics & Comp. Sci. & 3 & null \\ \hline
        CS-347 & null & null & null & CS-101 \\ \hline
        \end{tabular}
        \end{adjustbox}
    \end{center}
\end{frame}

\begin{frame}{Joined Relations}
    \begin{itemize}
        \item Join operations take two relations and return as a result another relation
        \item These additional operations are typically used as subquery expressions in the \texttt{from} clause
        \item \textbf{Join condition} - defines which tuples in the two relations match, and what attributes are present in the result of the join
        \item \textbf{Join type} - defines how tuples in each relation that do not match any tuple in the other relation (based on the join condition) are treated
    \end{itemize}
    
    \vspace{1em}
    \begin{columns}
        \begin{column}{0.5\textwidth}
            \textbf{Join Types}
            \begin{itemize}
                \item \texttt{inner join}
                \item \texttt{left outer join}
                \item \texttt{right outer join}
                \item \texttt{full outer join}
            \end{itemize}
        \end{column}
        \begin{column}{0.5\textwidth}
            \textbf{Join Conditions}
            \begin{itemize}
                \item \texttt{natural}
                \item \texttt{on <predicate>}
                \item \texttt{using (A1, A2, ..., An)}
            \end{itemize}
        \end{column}
    \end{columns}
\end{frame}

\begin{frame}[fragile]{Joined Relations - Examples}
    \begin{itemize}
        \item \texttt{course inner join prereq on course.course\_id = prereq.course\_id}
        \item What is the difference between the above (equi join), and a natural join?
    \end{itemize}
    \begin{center}
        \begin{adjustbox}{max width=\textwidth}
        \begin{tabular}{|l|l|l|l|l|l|}
        \hline
        \textbf{course\_id} & \textbf{title} & \textbf{dept\_name} & \textbf{credits} & \textbf{prereq\_id} & \textbf{course\_id} \\ \hline
        BIO-301 & Genetics & Biology & 4 & BIO-101 & BIO-301 \\ \hline
        CS-190 & Game Design & Comp. Sci. & 4 & CS-101 & CS-190 \\ \hline
        \end{tabular}
        \end{adjustbox}
    \end{center}

    \vspace{0.5em}
    \begin{itemize}
        \item \texttt{course left outer join prereq on course.course\_id = prereq.course\_id}
    \end{itemize}
    \begin{center}
        \begin{adjustbox}{max width=\textwidth}
        \begin{tabular}{|l|l|l|l|l|l|}
        \hline
        \textbf{course\_id} & \textbf{title} & \textbf{dept\_name} & \textbf{credits} & \textbf{prereq\_id} & \textbf{course\_id} \\ \hline
        BIO-301 & Genetics & Biology & 4 & BIO-101 & BIO-301 \\ \hline
        CS-190 & Game Design & Comp. Sci. & 4 & CS-101 & CS-190 \\ \hline
        CS-315 & Robotics & Comp. Sci. & 3 & null & null \\ \hline
        \end{tabular}
        \end{adjustbox}
    \end{center}
\end{frame}

\begin{frame}[fragile]{Joined Relations - Examples (2)}
    \begin{itemize}
        \item \texttt{course natural right outer join prereq}
    \end{itemize}
    \begin{center}
        \begin{adjustbox}{max width=0.9\textwidth}
        \begin{tabular}{|l|l|l|l|l|}
        \hline
        \textbf{course\_id} & \textbf{title} & \textbf{dept\_name} & \textbf{credits} & \textbf{prereq\_id} \\ \hline
        BIO-301 & Genetics & Biology & 4 & BIO-101 \\ \hline
        CS-190 & Game Design & Comp. Sci. & 4 & CS-101 \\ \hline
        CS-347 & null & null & null & CS-101 \\ \hline
        \end{tabular}
        \end{adjustbox}
    \end{center}

    \vspace{0.5em}
    \begin{itemize}
        \item \texttt{course full outer join prereq using (course\_id)}
    \end{itemize}
    \begin{center}
        \begin{adjustbox}{max width=0.9\textwidth}
        \begin{tabular}{|l|l|l|l|l|}
        \hline
        \textbf{course\_id} & \textbf{title} & \textbf{dept\_name} & \textbf{credits} & \textbf{prereq\_id} \\ \hline
        BIO-301 & Genetics & Biology & 4 & BIO-101 \\ \hline
        CS-190 & Game Design & Comp. Sci. & 4 & CS-101 \\ \hline
        CS-315 & Robotics & Comp. Sci. & 3 & null \\ \hline
        CS-347 & null & null & null & CS-101 \\ \hline
        \end{tabular}
        \end{adjustbox}
    \end{center}
\end{frame}

\section{Views}
\begin{frame}
  \vfill
  \centering
  \begin{beamercolorbox}[sep=8pt,center,shadow=true,rounded=true]{title}
    \usebeamerfont{title}Views\par%
  \end{beamercolorbox}
  \vfill
\end{frame}

\begin{frame}[fragile]{Views}
    \begin{itemize}
        \item In some cases, it is not desirable for all users to see the entire logical model (that is, all the actual relations stored in the database.)
        \item Consider a person who needs to know an instructor's name and department, but not the salary. This person should see a relation described, in SQL, by
        \begin{lstlisting}[language=SQL]
select ID, name, dept_name 
from instructor;
        \end{lstlisting}
        \item A \textbf{view} provides a mechanism to hide certain data from the view of certain users
        \item Any relation that is not of the conceptual model but is made visible to a user as a 'virtual relation' is called a \textbf{view}.
    \end{itemize}
\end{frame}

\begin{frame}[fragile]{View Definition}
    \begin{itemize}
        \item A view is defined using the \texttt{create view} statement which has the form
        \begin{lstlisting}[language=SQL]
create view v as <query expression>;
        \end{lstlisting}
        where \texttt{<query expression>} is any legal SQL expression
        \item The view name is represented by \texttt{v}
        \item Once a view is defined, the view name can be used to refer to the virtual relation that the view generates
        \item View definition is not the same as creating a new relation by evaluating the query expression
        \item Rather, a view definition causes the saving of an expression; the expression is substituted into queries using the view
    \end{itemize}
\end{frame}

\begin{frame}[fragile]{Example Views}
    \begin{itemize}
        \item A view of instructors without their salary
        \begin{lstlisting}[language=SQL]
create view faculty as 
select ID, name, dept_name 
from instructor;
        \end{lstlisting}
        \item Find all instructors in the Biology department
        \begin{lstlisting}[language=SQL]
select name 
from faculty 
where dept_name = 'Biology';
        \end{lstlisting}
        \item Create a view of department salary totals
        \begin{lstlisting}[language=SQL]
create view departments_total_salary(dept_name, total_salary) as 
select dept_name, sum(salary) 
from instructor 
group by dept_name;
        \end{lstlisting}
    \end{itemize}
\end{frame}

\begin{frame}[fragile]{Views Defined Using Other Views}
    \begin{lstlisting}[language=SQL, basicstyle=\ttfamily\scriptsize]
create view physics_fall_2009 as
select course.course_id, sec_id, building, room_number 
from course, section 
where course.course_id = section.course_id 
  and course.dept_name = 'Physics' 
  and section.semester = 'Fall' 
  and section.year = '2009';
    \end{lstlisting}
    
    \begin{lstlisting}[language=SQL, basicstyle=\ttfamily\scriptsize]
create view physics_fall_2009_watson as 
select course_id, room_number 
from physics_fall_2009 
where building = 'Watson';
    \end{lstlisting}
\end{frame}

\begin{frame}[fragile]{Views Defined Using Other Views (2)}
    \begin{itemize}
        \item One view may be used in the expression defining another view
        \item A view relation $v_1$ is said to depend directly on a view relation $v_2$ if $v_2$ is used in the expression defining $v_1$
        \item A view relation $v_1$ is said to depend on view relation $v_2$ if either $v_1$ depends directly on $v_2$ or there is a path of dependencies from $v_1$ to $v_2$
        \item A view relation $v$ is said to be \textbf{recursive} if it depends on itself
    \end{itemize}
\end{frame}

\begin{frame}[fragile]{View Expansion}
    \begin{itemize}
        \item Expand use of a view in a query/another view
        \begin{lstlisting}[language=SQL, basicstyle=\ttfamily\scriptsize]
create view physics_fall_2009_watson as 
( 
  select course_id, room_number 
  from (
      select course.course_id, building, room_number 
      from course, section 
      where course.course_id = section.course_id 
        and course.dept_name = 'Physics' 
        and section.semester = 'Fall' 
        and section.year = '2009'
  ) 
  where building = 'Watson'
);
        \end{lstlisting}
    \end{itemize}
\end{frame}

\begin{frame}{View Expansion*}
    \begin{itemize}
        \item A way to define the meaning of views defined in terms of other views
        \item Let view $v_1$ be defined by an expression $e_1$ that may itself contain uses of view relations
        \item View expansion of an expression repeats the following replacement step:
        \vspace{1em}
        \textbf{repeat}\\
        \quad Find any view relation $v_i$ in $e_1$\\
        \quad Replace the view relation $v_i$ by the expression defining $v_i$\\
        \textbf{until} no more view relations are present in $e_1$
        \vspace{1em}
        \item As long as the view definitions are not recursive, this loop will terminate
    \end{itemize}
\end{frame}

\begin{frame}[fragile]{Update of a View}
    \begin{itemize}
        \item Add a new tuple to \texttt{faculty} view which we defined earlier
        \begin{lstlisting}[language=SQL]
insert into faculty 
values ('30765', 'Green', 'Music');
        \end{lstlisting}
        \item This insertion must be represented by the insertion of the tuple \texttt{('30765', 'Green', 'Music', null)} into the \texttt{instructor} relation
    \end{itemize}
\end{frame}

\begin{frame}[fragile]{Some Updates cannot be Translated Uniquely}
    \begin{itemize}
        \item \begin{lstlisting}[language=SQL, basicstyle=\ttfamily\tiny]
create view instructor_info as 
select ID, name, building 
from instructor, department 
where instructor.dept_name = department.dept_name;
        \end{lstlisting}
        \item \texttt{insert into instructor\_info values ('69987', 'White', 'Taylor');}
        \item which department, if multiple departments in Taylor?
        \item what if no department is in Taylor?
        \item Most SQL implementations allow updates only on \textbf{simple views}
        \begin{itemize}
            \item The \texttt{from} clause has only one database relation
            \item The \texttt{select} clause contains only attribute names of the relation, and does not have any expressions, aggregates, or \texttt{distinct} specification
            \item Any attribute not listed in the \texttt{select} clause can be set to null
            \item The query does not have a \texttt{group by} or \texttt{having} clause
        \end{itemize}
    \end{itemize}
\end{frame}

\begin{frame}[fragile]{And Some Not at All}
    \begin{itemize}
        \item \begin{lstlisting}[language=SQL]
create view history_instructors as 
select * 
from instructor 
where dept_name = 'History';
        \end{lstlisting}
        \item What happens if we insert \texttt{('25566', 'Brown', 'Biology', 100000)} into \texttt{history\_instructors}?
    \end{itemize}
\end{frame}

\begin{frame}{Materialized Views}
    \begin{itemize}
        \item \textbf{Materializing a view}: create a physical table containing all the tuples in the result of the query defining the view
        \item If relations used in the query are updated, the materialized view result becomes out of date
        \item Need to \textbf{maintain} the view, by updating the view whenever the underlying relations are updated
    \end{itemize}
\end{frame}

\begin{frame}{Module Summary}
    \begin{itemize}
        \item Learnt SQL expressions for Join and Views
    \end{itemize}
    \vspace{2em}
    \begin{center}
    \small \textit{Slides used in this presentation are borrowed from \url{http://db-book.com/} with kind permission of the authors.} \\
    \small \textit{Edited and new slides are marked with 'PPD'.}
    \end{center}
\end{frame}

\end{document}
